\section{Organisation}

\begin{frame}{Disclaimer - Quellen dieser Unterlagen}
  \begin{alertblock}{Wichtiger Hinweis zu den Inhalten}
    \begin{itemize}
      \item Diese Materialien basieren direkt auf den Vorlesungsunterlagen von Markus Brenneis
      \item Die Folien basieren auf den Folien von Sven Hoops und Jonas Goudarzi
    \end{itemize}
  \end{alertblock}
\end{frame}
  
\begin{frame}{Wie für die Klausur lernen (offizielle Version)}
  \begin{alertblock}{Nur Altklausuren lernen = Schiefgehen!}
    \begin{itemize}
      \item Altklausuren sind \textbf{kein Ersatz} für intensive Mitarbeit im Semester
      \item Klausuren prüfen \textbf{tieferes Verständnis} ab, nicht nur Auswendiglernen
      \item \textbf{Alle Wochenblätter sind relevant}, auch die letzten Blätter
    \end{itemize}
  \end{alertblock}

  \begin{exampleblock}{Richtige Lernstrategie}
    \begin{itemize}
      \item Von jedem Woche die \textbf{Wichtigen Begriffe} erklären können \& die \textbf{Lernziele} beachten
      \item \textbf{Alle Aufgaben selbstständig} bearbeiten
      \item \textbf{Anderen den Stoff erklären} (Lerngruppen!)
      \item Muster in den Altklausuren erkennen und verstehen, wie man die Aufgaben angeht
    \end{itemize}
  \end{exampleblock}
\end{frame}

\begin{frame}{Nachklausurtutorium Programmierung}
  \textbf{Zeitraum:} Montag den 09.03.2026 bis Freitag den 13.03.2026

  \textbf{Zeit:} Täglich 10:30 bis 14:00 Uhr

  \textbf{Ort:} Hörsaal 5H

  \vspace{0.5cm}

\end{frame}